% vip-ribs-mvi-redux.tex — map of iOS architectures beyond MVC/MVVM/VIPER: Clean Swift (VIP),
% RIBs (Uber), MVI, Redux/ReSwift, TCA and the SwiftUI "MV" approach. Four loops side by side,
% a who-owns-state table, and the lineage of the ideas.
% Sources (checked 2026-10-05):
%   clean-swift.com/clean-swift-ios-architecture · github.com/Clean-Swift/CleanStore (ListOrders scene)
%   github.com/uber/RIBs (README, wiki) · github.com/uber/RIBs-iOS RIBs/Classes/Interactor.swift
%   uber.com/blog/new-rider-app-architecture (2016-12-20)
%   staltz.com/unidirectional-user-interface-architectures.html (2015-08-22)
%   hannesdorfmann.com/android/mosby3-mvi-1 · redux.js.org (three principles; history of Redux)
%   github.com/ReSwift/ReSwift (README, releases) · github.com/pointfreeco/swift-composable-architecture
%   developer.apple.com/documentation/swiftui/managing-model-data-in-your-app
%   objc.io/issues/13-architecture/viper (Gilbert, Stoll, 2014-06) · redux.js.org/faq/organizing-state
%   developer.android.com/topic/architecture/ui-layer/events
% @labels: area=design kind=architecture platform=ios discussion=5
% @tags: vip, clean-swift, ribs, mvi, redux, reswift, mv-pattern, viper, interactor, unidirectional-data-flow, tca, observation
\documentclass[8pt]{extarticle}
\usepackage{printup-sheet}
\usepackage{array}

\lstdefinelanguage{SwiftSheet}{
  morekeywords={protocol,class,final,struct,enum,func,var,let,weak,init,
    if,else,return,guard,self,nil,try,await,async,private,some,is,
    true,false,case,switch,default,static,in,AnyObject},
  sensitive=true, morecomment=[l]{//}, morestring=[b]",
  literate={->}{{\hbox{-}\hbox{>}}}2 {==}{{\hbox{=}\hbox{=}}}2 {??}{{\hbox{?}\hbox{?}}}2}

\tikzset{
  rb/.style={box, font=\footnotesize, minimum width=15mm, minimum height=6.5mm, inner sep=2pt},
  vw/.style={rb, draw=sheetBlue, fill=sheetBlue!8},
  lg/.style={rb, draw=sheetGreen!60!black, fill=sheetGreen!12},
  pr/.style={rb, draw=sheetOrange, fill=sheetOrange!10},
  rt/.style={rb, draw=sheetBrown, fill=sheetBrown!10},
  lbl/.style={font=\scriptsize, text=black!75, inner sep=1pt, align=center},
  ttl/.style={font=\bfseries\normalsize},
  tn/.style={draw=sheetGreen!60!black, fill=sheetGreen!10, rounded corners=2pt, font=\scriptsize,
             inner sep=1.5pt, align=center, minimum width=12mm, minimum height=5.5mm},
  ev/.style={circle, fill=sheetBlue, inner sep=1.6pt},
  yr/.style={font=\bfseries\scriptsize, text=sheetBlue},
  evl/.style={font=\scriptsize, align=center, text=black!85},
}

\begin{document}

\sheettitle{VIP · RIBs · MVI · Redux · MV — the architecture map}{design · folio}

\oneliner{Every iOS architecture answers three questions: \textbf{who owns the state},
\textbf{which way does data flow}, and \textbf{where do side effects live}. \textbf{VIP} (Clean
Swift) makes each screen a one-way \emph{cycle} of three objects; \textbf{RIBs} (Uber) drive the app
by a \emph{tree of business logic}, not of views; \textbf{MVI} and \textbf{Redux} keep one
immutable state changed only by a reducer; \textbf{TCA} is Redux made composable; \textbf{MV} says
SwiftUI views plus \texttt{@Observable} models are already enough.}

\vspace{2pt}
\noindent\begin{tikzpicture}[sheet, yscale=1.12]
  % ── 1. VIP ─────────────────────────────────────────────
  \node[ttl] at (1.9,2.75) {VIP (Clean Swift)};
  \node[vw] (vc) at (1.9,2.05) {ViewController};
  \node[lg] (in) at (3.35,0.75) {Interactor};
  \node[pr] (pe) at (0.45,0.75) {Presenter};
  \draw[flow] (vc.east) -| node[lbl, right, pos=0.7]{Request} (in.north);
  \draw[flow] (in) -- node[lbl, below]{Response} (pe);
  \draw[hot] (pe.north) |- node[lbl, left, pos=0.3]{ViewModel} node[lbl, above, pos=0.75, text=sheetOrange]{weak} (vc.west);
  \node[rt, minimum width=11mm] (ro) at (1.2,-0.4) {Router\\[-2pt]{\scriptsize owned by VC}};
  \node[rb, minimum width=11mm, draw=sheetGrey, fill=black!4] (wo) at (3.35,-0.4) {Worker\\[-2pt]{\scriptsize API, DB}};
  \draw[flow] (in) -- (wo);
  \node[lbl, text=sheetGreen!50!black] at (1.9,-1.05) {one-way cycle, one per scene};

  % ── 2. RIBs ────────────────────────────────────────────
  \node[ttl] at (6.4,2.75) {RIBs (Uber)};
  \node[tn] (root) at (6.4,2.05) {Root\\R · I · V};
  \node[tn] (lo) at (5.2,1.05) {LoggedOut\\R · I · V};
  \node[tn, draw=sheetOrange, fill=sheetOrange!10, dashed] (li) at (7.6,1.05) {LoggedIn\\R · I \textbf{(no view)}};
  \node[tn] (og) at (6.9,0.0) {OffGame\\R · I · V};
  \node[tn] (tt) at (8.3,0.0) {Game\\R · I · V};
  \draw[flow] (root) -- (lo); \draw[flow] (root) -- (li);
  \draw[flow] (li) -- (og); \draw[flow] (li) -- (tt);
  \draw[hot] (tt.north) to[bend right=40] (li.east);
  \node[lbl, align=center] at (5.0,-0.05) {Router\\attaches /\\detaches\\children};
  \node[lbl, text=sheetGreen!50!black] at (6.5,-0.78) {tree of \emph{logic}; views follow};
  \node[lbl] at (6.5,-1.07) {Builder = factory + DI · \textcolor{sheetOrange}{orange = Listener (up)}};

  % ── 3. MVI ─────────────────────────────────────────────
  \node[ttl] at (10.7,2.75) {MVI};
  \node[vw] (mv) at (10.7,2.05) {View\\[-1pt]{\scriptsize \texttt{render(state)}}};
  \node[pr] (mi) at (12.0,0.65) {Intent\\[-1pt]{\scriptsize events → actions}};
  \node[lg] (mm) at (9.5,0.65) {Model\\[-1pt]{\scriptsize actions → state}};
  \draw[flow] (mv.east) -| node[lbl, right, pos=0.7]{events} (mi.north);
  \draw[flow] (mi) -- node[lbl, above]{actions} (mm);
  \draw[flow] (mm.north) |- node[lbl, left, pos=0.3]{state} (mv.west);
  \node[lbl, text=sheetGreen!50!black] at (10.7,-0.78) {one immutable state per screen};
  \node[lbl] at (10.7,-1.07) {RxJS Observables in Cycle.js};

  % ── 4. Redux / ReSwift ────────────────────────────────
  \node[ttl] at (14.6,2.75) {Redux / ReSwift};
  \node[vw] (rv) at (14.6,2.05) {View\\[-1pt]{\scriptsize \texttt{StoreSubscriber}}};
  \node[rb, draw=sheetGrey, fill=black!4] (mw) at (15.85,0.95) {Middleware};
  \node[lg] (rr) at (15.85,-0.15) {Reducer\\[-1pt]{\scriptsize pure}};
  \node[rt] (rs) at (13.65,0.35) {Store\\[-1pt]{\scriptsize \textbf{one} app State}};
  \draw[flow] (rv.east) -| node[lbl, left, pos=0.93]{\texttt{dispatch(a)}} (mw.north);
  \draw[flow] (mw) -- (rr);
  \draw[flow] (rr.west) -- (rs.south east);
  \draw[flow] (rs.north) |- node[lbl, left, pos=0.3, align=right]{\texttt{newState}\\\texttt{(state:)}} (rv.west);
  \node[lbl, text=sheetGreen!50!black] at (14.6,-0.78) {single source of truth};
  \node[lbl] at (14.6,-1.07) {async: thunks (ReSwift-Thunk)};
\end{tikzpicture}

\vspace{1pt}
{\footnotesize\setlength{\tabcolsep}{3pt}\renewcommand{\arraystretch}{1.05}
\noindent\begin{tabular}{@{}>{\raggedright\arraybackslash}p{17mm}>{\raggedright\arraybackslash}p{27mm}>{\raggedright\arraybackslash}p{30mm}>{\raggedright\arraybackslash}p{29mm}>{\raggedright\arraybackslash}p{21mm}>{\raggedright\arraybackslash}p{35mm}@{}}
\toprule
\textbf{Pattern} & \textbf{State lives in} & \textbf{Data flow} & \textbf{Testability} & \textbf{Boilerplate} & \textbf{Fits} \\
\midrule
VIP & Interactor (+ \texttt{DataStore}) & VC$\to$I$\to$P$\to$VC, one way & high: spy the next role & high: 7 per scene & UIKit apps, per-scene rigour \\
RIBs & each RIB's Interactor & tree; Rx down, listener up & high: I/R/B behind protocols & very high, codegen & huge apps, hundreds of engineers \\
MVI & one immutable screen State & Intent$\to$Model$\to$View & high: pure reducer & medium & complex screens, Android parity \\
Redux / ReSwift & one global Store & dispatch$\to$reducer$\to$subscribers & high reducers; effects harder & medium–high & shared app-wide state \\
TCA & Store tree of features & unidirectional, composed & very high: \texttt{TestStore} & high & big SwiftUI teams \\
MV & \texttt{@Observable} models & view calls model, observes it & models yes, view logic no & low & SwiftUI small–mid apps \\
\bottomrule
\end{tabular}}

\vspace{2pt}
\noindent\begin{tikzpicture}[sheet]
  \draw[thick, black!40] (-0.2,0) -- (16.6,0);
  \draw[black!40, thick] (1.25,-0.13) -- (1.45,0.13);
  \draw[black!40, thick] (1.45,-0.13) -- (1.65,0.13);
  \node[ev] at (0.4,0) {};   \node[yr, above] at (0.4,0.12) {1979};
  \node[evl, below] at (0.4,-0.12) {\textbf{MVC} · Reenskaug\\the View observes the Model};
  \node[ev] at (3.3,0) {};   \node[yr, above] at (3.3,0.12) {2014};
  \node[evl, below] at (3.3,-0.12) {\textbf{Flux} · Facebook\\one-way flow};
  \node[ev] at (6.1,0) {};   \node[yr, above] at (6.1,0.12) {mid-2015};
  \node[evl, below] at (6.1,-0.12) {\textbf{Redux} · Abramov\\one store, pure reducers};
  \node[ev] at (8.9,0) {};   \node[yr, above] at (8.9,0.12) {Aug 2015};
  \node[evl, below] at (8.9,-0.12) {\textbf{MVI} · Staltz\\Cycle.js, RxJS};
  \node[ev] at (11.7,0) {};  \node[yr, above] at (11.7,0.12) {Dec 2016};
  \node[evl, below] at (11.7,-0.12) {\textbf{RIBs} · Uber\\rider-app rewrite};
  \node[ev] at (14.6,0) {};  \node[yr, above] at (14.6,0.12) {iOS 17};
  \node[evl, below] at (14.6,-0.12) {\textbf{\texttt{@Observable}}\\SwiftUI ``MV'', no ViewModel};
  \node[font=\tiny, text=black!50, anchor=east] at (16.6,0.32) {not to scale};
\end{tikzpicture}

\begin{multicols}{2}
\raggedright

\section{How it works}
\begin{itemize}
  \item \textbf{VIP} (Clean Swift; Clean Architecture per scene): VC
        (\texttt{DisplayLogic}) sends a \emph{Request} to the Interactor (\texttt{BusinessLogic},
        uses \textbf{Workers}, holds the \texttt{DataStore}); its \emph{Response} goes to the
        Presenter (\texttt{PresentationLogic}), which formats a \emph{ViewModel} and calls
        \texttt{display…} on the VC. \textbf{Router} (\texttt{RoutingLogic} +
        \texttt{DataPassing}) navigates, passing data via the \texttt{DataStore}. Models nest per
        use case (\texttt{ListOrders.FetchOrders.Request}). Older templates: \texttt{…Input}/%
        \texttt{…Output} protocols, wired by a \textbf{Configurator}.
  \item \textbf{RIBs} = \textbf{R}outer · \textbf{I}nteractor · \textbf{B}uilder (+ Component for
        dependencies, optional Presenter and View). The Interactor holds business logic and Rx
        subscriptions, alive between \texttt{didBecomeActive()} and \texttt{willResignActive()};
        bind with \texttt{disposeOnDeactivate(interactor:)}. Router attaches/detaches children;
        Builder creates the unit and is the only part that knows the DI system. Up: a
        \texttt{Listener} the parent implements; down: Rx streams or \texttt{build()} arguments.
  \item \textbf{MVI}: \emph{Intent} turns user events into actions, \emph{Model} turns actions
        into state, \emph{View} renders state — a pure \texttt{render(state)} of one immutable
        state.
  \item \textbf{Redux}'s 3 principles: one store = single source of truth · state is
        read-only, changed only by an action · pure reducers (state + action $\to$ new
        state). ReSwift: \texttt{Store<AppState>(reducer:state:)}, \texttt{dispatch},
        \texttt{newState(state:)}; \texttt{subscribe(self) \{ \$0.select(…) \}} narrows
        updates, skipping repeats of an \texttt{Equatable} substate.
  \item \textbf{MV}: no per-screen ViewModel; logic lives in \texttt{@Observable} models
        (iOS 17, macOS 14). A view updates only when a property its \texttt{body} reads
        changes. Own with \texttt{@State}, share with \texttt{.environment(model)}, read with
        \texttt{@Environment(Model.self)}, bind with \texttt{@Bindable}.
  \item \textbf{TCA}: State · Action · Reducer · Store; effects and dependencies are
        injectable, and \texttt{TestStore} asserts every state change and effect.
\end{itemize}

\section{Why RIBs exist}
Uber's rider app (rewrite, 2016): an MVC app grown to a few hundred engineers, one
\texttt{RequestViewController} from 300 to 3,000+ lines. RIBs make \textbf{business logic, not
the view tree, drive routing} — a deep logic tree with a shallow view tree.

\section{Example — one VIP use case}
\begin{lstlisting}[language=SwiftSheet]
enum ListOrders { enum Fetch {          // per use case
  struct Request {}
  struct Response { let orders: [Order] }
  struct ViewModel { let rows: [String] } } }
protocol ListOrdersDisplayLogic: AnyObject {
  func display(_ vm: ListOrders.Fetch.ViewModel) }
final class ListOrdersInteractor {      // BusinessLogic
  var presenter: ListOrdersPresenter?; let worker = OrdersWorker()
  func fetch(_ r: ListOrders.Fetch.Request) {
    worker.load { self.presenter?.present(.init(orders: $0)) } } }
final class ListOrdersPresenter {       // PresentationLogic
  weak var viewController: ListOrdersDisplayLogic?  // no cycle
  func present(_ r: ListOrders.Fetch.Response) {
    viewController?.display(.init(rows: r.orders.map(\.title))) } }
\end{lstlisting}

\section{Traps}
\begin{itemize}
  \trap{``VIP is VIPER renamed.'' VIPER (2014): Presenter \emph{between} View and Interactor,
        both ways; Entities stay in the Interactor. VIP: a one-way cycle.}
  \trap{RIBs' Router is not ``just navigation'': it attaches/detaches child RIBs from
        Interactor output — screens follow the tree.}
  \trap{Redux: subscribers hear every change unless they \texttt{select} a slice; local UI
        state (\texttt{isDropdownOpen}) stays out of the store.}
  \trap{A one-off message kept in state shows again on every re-render until the view
        reports it shown and the state clears it.}
  \trap{Conforming to the \texttt{Observable} protocol alone adds no observation — use the
        \texttt{@Observable} macro.}
\end{itemize}

\end{multicols}

\noindent{\footnotesize\color{sheetGrey}\textit{Related:} MVC / MVP / MVVM · Coordinator, Repository,
DI, Clean / VIPER · The Composable Architecture (TCA) · unidirectional state machines · Modularization \& SPM}

\end{document}
